%% Question bank chapter file - generated by bank_tools.py build (do not edit while the bank is still being built:
%% it is regenerated from the per-test data in Question Bank\data). Every question is one bankq block:
%%   \begin{bankq}[id=..., source={...}, paper=.., qno=.., type=..., subtopic={...}, answer={...}, answer_src=key|solved,
%%                 working={...}, difficulty=1-10, diff_parts={K S I C F}, diff_note={...},
%%                 group=<passage id or empty>, also_in={duplicate copies in other tests}]
%% difficulty = min(10, K+S+I+C+F-1): K concepts 1-3, S steps 1-4, I insight 0-2, C computation 0-2, F format 0-1
%%   ...question LaTeX, options, figures (\bfig / \bfigi -> img/)...
%%   \end{bankq}
%% Passages shared by several questions are bankpassage blocks; their questions carry group=<passage id>.
\documentclass[10pt]{article}
\usepackage[a4paper,margin=1.8cm]{geometry}
\usepackage[T1]{fontenc}
\usepackage{amsmath,amssymb,mathtools}
\usepackage{graphicx}
\usepackage[export]{adjustbox}
\usepackage{xcolor,needspace,keyval}
\usepackage{enumitem}
\usepackage{booktabs,array,tabularx}
\usepackage[hidelinks]{hyperref}
\setlength{\parindent}{0pt}
\setlength{\parskip}{3pt}
\makeatletter
\define@key{bq}{id}{\def\bq@id{#1}}\define@key{bq}{source}{\def\bq@src{#1}}\define@key{bq}{paper}{\def\bq@paper{#1}}
\define@key{bq}{qno}{\def\bq@qno{#1}}\define@key{bq}{type}{\def\bq@type{#1}}\define@key{bq}{subtopic}{\def\bq@sub{#1}}
\define@key{bq}{answer}{\def\bq@ans{#1}}\define@key{bq}{answer_src}{\def\bq@asrc{#1}}\define@key{bq}{working}{\def\bq@work{#1}}
\define@key{bq}{group}{\def\bq@group{#1}}\define@key{bq}{also_in}{\def\bq@also{#1}}
\define@key{bq}{difficulty}{\def\bq@diff{#1}}\define@key{bq}{diff_parts}{\def\bq@dparts{#1}}\define@key{bq}{diff_note}{\def\bq@dnote{#1}}
\newenvironment{bankq}[1][]{\def\bq@work{}\def\bq@also{}\def\bq@group{}\def\bq@diff{?}\def\bq@dparts{}\def\bq@dnote{}\setkeys{bq}{#1}%
  \par\needspace{4\baselineskip}\medskip\noindent{\small\textbf{\bq@id}\hfill\color{gray}[\bq@src, Paper \bq@paper, Q\bq@qno]
  \bq@type\ $\cdot$ \bq@sub\ $\cdot$ \textbf{D\bq@diff}}\par\nopagebreak\relax}
  {\par{\small\color{teal}Answer: \textbf{\bq@ans}\ifx\bq@work\@empty\else\ (solved: \bq@work)\fi
  \ifx\bq@also\@empty\else\par\color{gray}Also in: \bq@also\fi}\par}
\define@key{bp}{id}{\def\bp@id{#1}}\define@key{bp}{source}{\def\bp@src{#1}}\define@key{bp}{questions}{\def\bp@qs{#1}}
\newenvironment{bankpassage}[1][]{\setkeys{bp}{#1}\par\needspace{4\baselineskip}\bigskip\noindent
  {\small\textbf{Passage \bp@id}\hfill\color{gray}[\bp@src; questions \bp@qs]}\par\nopagebreak\relax\itshape}{\par}
\makeatother
\newcommand{\bfig}[2]{\par\smallskip{\centering\includegraphics[width=#1pt,max width=0.62\linewidth,max height=0.32\textheight]{img/#2}\par}\smallskip}
\newcommand{\bfigi}[2]{\includegraphics[width=#1pt,max width=\linewidth,max height=0.32\textheight,valign=c]{img/#2}}
\newcommand{\bmissing}[1]{\textit{[figure missing in the source conversion: \detokenize{#1}]}}
\newenvironment{qoptions}{\begin{itemize}[leftmargin=1.8em,itemsep=1pt,topsep=2pt,label={}]}{\end{itemize}}
\newcommand{\opt}[2]{\item[(#1)] #2}
\providecommand{\fwstart}{}\providecommand{\fwend}{}\providecommand{\qfullsplit}{}
\newcolumntype{L}[1]{>{\raggedright\arraybackslash}p{#1}}
\newcolumntype{C}[1]{>{\centering\arraybackslash}p{#1}}
\newcolumntype{R}[1]{>{\raggedleft\arraybackslash}p{#1}}
\providecommand{\sectionhdr}[1]{}\providecommand{\subsectionhdr}[1]{}


\begin{document}
{\Large\bfseries Heat Transfer --- Question Bank}\par\smallskip
{\small Physics chapter P16. Questions from all converted test papers, grouped by source series. Answers: printed key where it matches the paper, otherwise solved (answer\_src=solved, with working). Exact repeats across papers appear once, with also\_in listing the other copies.}\par\medskip\hrule

\section*{Score Advanced papers}
\begin{bankq}[id=SA02-P1-PHY-S1-Q04, source={Score Advanced - Test 02 (13 May 2021)}, paper=1, qno=4, type=multi, subtopic={Conduction}, answer={A,C,D}, answer_src=solved, working={Heat into A: 1+3+1(from E)=5 out of D; flows are proportional to temperature differences: TB=TE>TA, TC>TD}, difficulty=5, diff_parts={K1 S2 I1 C1 F1}, diff_note={concepts: heat flow through rods meeting at a junction | steps: heat current proportional to temperature difference; balance at the junction; compare temperatures}]
Four identical rods which have thermally insulated lateral surfaces are joined at point A. Points B,C,D \& E are connected to large reservoirs. If heat flows into the junction from point B at rate of 1 W and from point C at 3 W inside, flows out from D at 5 W, which relation(s) is/are correct for temperature of these points ?\par
\bfig{56.6}{SA02_p1_phy_s1_q04_1.png}
\begin{qoptions}
\opt{A}{T$_\mathrm A<$ T$_\mathrm E$}
\opt{B}{T$_\mathrm B$ = T$_\mathrm C$}
\opt{C}{T$_\mathrm C>$ T$_\mathrm D$}
\opt{D}{T$_\mathrm B$ = T$_\mathrm E$}
\end{qoptions}
\end{bankq}
\begin{bankq}[id=SA05-P1-PHY-S1-Q14, source={Score Advanced - Test 05 (15 Jul 2021)}, paper=1, qno=14, type=single, subtopic={Conduction}, answer={C}, answer_src=key, difficulty=4, diff_parts={K2 S1 I1 C1 F0}, diff_note={concepts: interface temperature of a composite rod, conduction | steps: heat current is equal in both sections}]
One end of a thermally insulated rod is kept at a temperature T$_1$ and the other at T$_2$. The rod is composed of two sections of lengths $\ell_1$ and $\ell_2$ and thermal conductivities K$_1$ and K$_2$ respectively. The temperature at the interface of the two sections is-\par
\bfig{126.3}{SA05_p1_phy_s1_q14_1.png}
\begin{qoptions}
\opt{A}{(K$_2\ell_2$T$_1$ + K$_1\ell_1$T$_2$)/(K$_1\ell_2$ + K$_2\ell_1$)}
\opt{B}{(K$_2\ell_1$T$_1$ + K$_1\ell_2$T$_2$)/(K$_2\ell_1$ + K$_1\ell_2$)}
\opt{C}{(K$_1\ell_2$T$_1$ + K$_2\ell_1$T$_2$)/(K$_1\ell_2$ + K$_2\ell_1$)}
\opt{D}{(K$_1\ell_1$T$_1$ + K$_2\ell_2$T$_2$)/(K$_1\ell_1$ + K$_2\ell_2$)}
\end{qoptions}
\end{bankq}
\begin{bankq}[id=SA08-P1-PHY-S3-Q01, source={Score Advanced - Test 08 (22 Apr 2023)}, paper=1, qno=1, type=numerical, subtopic={Conduction}, answer={7}, answer_src=key, difficulty=9, diff_parts={K3 S2 I2 C2 F1}, diff_note={concepts: temperature difference inside a radioactive planet, conduction, heat flow through a spherical shell | steps: heat flow through a spherical shell; integrate the conduction equation; surface temperature from radiation}]
A certain planet of radius R is composed of a uniform material that, through radioactive decay, generates a net power P such that heat per second $\dfrac{\Delta\text{Q}}{\Delta\text{t}}$ = P$\dfrac{\text{r}^3}{\text{R}^3}$, where r is radial distance from centre. This results in a temperature differential between the inside and outside of the planet as heat is transferred from the interior to the surface. The rate of heat transfer inside is governed by conduction. It is found that thermal conductivity k is constant for the planet. For the following assume that the planet is in a steady state; temperature might depend on position, but does not depend on time. (Assuming black body radiation, and emissivity is 1) Find an expression for the temperature difference between the surface of the planet (T$_\text{S}$) and the centre of the planet (T$_\text{C}$). If your answer is $\Delta$T = $\dfrac{\sigma\text{R}^\text{m}\text{T}_\text{S}^\text{n}}{\ell\text{k}}$ fill value of (m + n + $\ell$)
\end{bankq}
\begin{bankq}[id=SA08-P2-PHY-S1-Q04, source={Score Advanced - Test 08 (22 Apr 2023)}, paper=2, qno=4, type=multi, subtopic={Conduction}, answer={A,B,C}, answer_src=key, difficulty=8, diff_parts={K3 S1 I2 C2 F1}, diff_note={concepts: two reservoirs joined by a rod, conduction, heat flow through the rod | steps: heat flow through the rod; thermal time constant}]
Figure shows two reservoirs containing two liquids of masses 4m and 2m and their specific heat capacities are S and 2S respectively. Their initial temperatures are 4T$_0$ and T$_0$ respectively. The containers are joined by a conducting rod of thermal conductivity K, length $\ell$ and cross-section area A. Specific heat capacity of the rod is negligible. (Neglect radiation)\par
\bfig{139.8}{SA08_p2_phy_s1_q04_1.png}
\begin{qoptions}
\opt{A}{The temperature of mid-point `P' of the rod is always constant.}
\opt{B}{Heat extracted from container 1 = Heat supplied to container 2.}
\opt{C}{After time t$_0$ = $\dfrac{2\text{ms}\ell}{\text{kA}}\ell$n2 the temperature difference between containers becomes half of the initial temperature difference.}
\opt{D}{After time t$_0$ = $\dfrac{\text{ms}\ell}{\text{kA}}\ell$n2 the temperature difference between containers becomes half of the initial temperature difference.}
\end{qoptions}
\end{bankq}
\begin{bankq}[id=SA09-P2-PHY-S1-Q06, source={Score Advanced - Test 09 (29 Apr 2023)}, paper=2, qno=6, type=multi, subtopic={Stefan-Boltzmann law}, answer={A,B,C}, answer_src=key, difficulty=8, diff_parts={K3 S1 I2 C2 F1}, diff_note={concepts: heat balance of a current carrying cylinder inside a hollow cylinder, stefan-boltzmann law, power radiated equals Joule heating | steps: power radiated equals Joule heating; view factors}]
A long uniform solid metallic cylinder has cross sectional radius r. Surface of the cylinder is painted by black body material. Resistivity of the material of inner cylinder is $\rho$. Another long hollow coaxial cylinder of black body material has cross sectional radius R = 8r. Electric current I continuously passes through the axis of inner cylinder. Whole arrangement lies in vacuum. When steady state is reached temperature of outer shell is 500 K. Choose the correct answer.\par
\bfig{147.9}{SA09_p2_phy_s1_q06_1.png}
\begin{qoptions}
\opt{A}{Temperature of surface of inner cylinder is 500$\sqrt{3}$ K.}
\opt{B}{Current through the inner cylinder is $\sqrt{\dfrac{\sigma}{\rho}16\pi^2\text{r}^3(500)^4}$.}
\opt{C}{Intensity of emitted radiation at distance 16r from axis is $\dfrac{\sigma}{2}(500)^4$.}
\opt{D}{None of these.}
\end{qoptions}
\end{bankq}
\begin{bankq}[id=SA12-P1-PHY-S1-Q04, source={Score Advanced - Test 12 (11 May 2023)}, paper=1, qno=4, type=single, subtopic={Stefan-Boltzmann law}, answer={B}, answer_src=key, difficulty=7, diff_parts={K3 S1 I2 C2 F0}, diff_note={concepts: bodies with different emissivities cooling by radiation, stefan-boltzmann law, rate of fall of temperature from Stefan's law | steps: rate of fall of temperature from Stefan's law}]
Two bodies A and B have the same surface area and mass. The bodies have absolute temperature T$_\text{A}$ and T$_\text{B}$ , emissivities e$_\text{A}$ and e$_\text{B}$ , and specific heat capacities S$_\text{A}$ and S$_\text{B}$ . The intensity of radiation (E) is plotted against the wavelength ($\lambda$) of radiation for both bodies. Which of the following is possible?\par
\bfig{79.3}{SA12_p1_phy_s1_q04_1.png}
\begin{qoptions}
\opt{A}{T$_\text{A}$ = T$_\text{B}$, e$_\text{A}$ = e$_\text{B}$}
\opt{B}{T$_\text{A}$ = T$_\text{B}$, e$_\text{A}$ $\neq$ e$_\text{B}$}
\opt{C}{T$_\text{A}$ $\neq$ T$_\text{B}$, e$_\text{A}$ = e$_\text{B}$,}
\opt{D}{T$_\text{A}$ $\neq$ T$_\text{B}$, e$_\text{A}$ = e$_\text{B}$}
\end{qoptions}
\end{bankq}
\begin{bankq}[id=SA13-P1-PHY-S2-Q02, source={Score Advanced - Test 13 (13 May 2023)}, paper=1, qno=2, type=numerical, subtopic={Conduction}, answer={1}, answer_src=key, difficulty=8, diff_parts={K3 S1 I2 C2 F1}, diff_note={concepts: heat conducted through two welded rods in series, conduction, equal thermal currents | steps: equal thermal currents; temperature drop at the joint}]
Two identical rectangular rods of metal are welded end to end in series between temperature 0$^\circ$C and 100$^\circ$C and 10 J of heat is conducted (in steady state process) through the rod in 2.00 min. If 5 such rods are taken and joined as shown in figure maintaining the same temperature difference between A and B, then the time (in min) in which 20 J heat will flow through the rods is :-\par
\bfig{68.8}{SA13_p1_phy_s2_q02_1.png}
\end{bankq}

\section*{TOAS papers}
\begin{bankq}[id=TOAS11-P1-PHY-S2-Q10, source={TOAS - Test 11 (27 Jun 2022)}, paper=1, qno=10, type=multi, subtopic={Conduction}, answer={A,C,D}, answer_src=key, difficulty=5, diff_parts={K2 S2 I1 C0 F1}, diff_note={concepts: conduction through a spherical ice shell, latent heat | steps: heat current through a spherical shell of ice; heat released by freezing a thickness dr; integrate to get the time; time for half thickness | I: heat conduction through a shell of changing thickness}]
A highly thin conducting spherical shell of radius $R_0$ is completely filled with water at $0^\circ$C. Outside surrounding temperature is $-T_0^\circ$C. Latent heat of fusion of water is $L_0$ and thermal conductivity of ice is $K_{\text{ice}}$. Assume heat is flowing symmetrically in radial direction. Choose the CORRECT option(s). Density of water is $\rho_0$. Neglect the expansion of water on freezing.
\bfig{258.6}{TOAS11_p1_phy_s2_q10_1.png}

\begin{qoptions}
\opt{A}{Time in which complete water will freeze, is given by $\dfrac{\rho_0L_0R_0^2}{6K_{\text{ice}}T_0}$}
\opt{B}{Time in which complete water will freeze, is given by $\dfrac{\rho_0L_0R_0^2}{4K_{\text{ice}}T_0}$}
\opt{C}{Time in which thickness of ice will become $\dfrac{R_0}{2}$, is given by $\dfrac{\rho_0L_0R_0^2}{12K_{\text{ice}}T_0}$}
\opt{D}{Total heat released by water in freezing, is given by $\dfrac{4\pi\rho_0L_0R_0^3}{3}$}
\end{qoptions}
\end{bankq}
\begin{bankq}[id=TOAS11-P1-PHY-S3-Q04, source={TOAS - Test 11 (27 Jun 2022)}, paper=1, qno=4, type=numerical, subtopic={Newton's law of cooling}, answer={2.00}, answer_src=key, difficulty=3, diff_parts={K1 S2 I0 C0 F1}, diff_note={concepts: heating with Newton's law of cooling | steps: C dT/dt = P - kT; T(t); time to half of the maximum temperature}]
A highly conducting uniform sphere of thermal capacity $C$ is heated by an electric heater, a resistance $R$ fitted within the sphere. A constant current $I$ is passed through the heater starting at time $t=0$ which gives constant power. The sphere loses heat at a rate equal to $k$ times the temperature difference between the sphere and the surrounding. The initial temperature of the sphere and that of the surrounding is $0^\circ$C. The time at which sphere attains half of its maximum attainable temperature (in $^\circ$C) is $\dfrac ck\ln a$. Then find the value of ``$a$''.
\end{bankq}
\begin{bankq}[id=TOAS11-P1-PHY-S3-Q05, source={TOAS - Test 11 (27 Jun 2022)}, paper=1, qno=5, type=numerical, subtopic={Stefan-Boltzmann law}, answer={1.60 to 1.70}, answer_src=key, difficulty=4, diff_parts={K2 S2 I0 C0 F1}, diff_note={concepts: net radiation absorbed by a body, heat capacity | steps: net radiative power from the enclosure; add the supplied power; divide by the heat capacity}]
A body of mass 2 gm has specific heat capacity 0.25 cal/gm-$^\circ$C and surface area 30 cm$^2$. It is kept in an enclosure which is maintained at temperature 127$^\circ$C. Body is being supplied heat at rate 2 watt. If temperature of body is 27$^\circ$C and emissivity of body is 0.5. Determine the rate at which its temperature is changing (in K/s) at the given instant. (Stefan's constant is $5.67\times10^{-8}$ W/m$^2$-K$^4$. 1 cal = 4.2 J)
\end{bankq}
\begin{bankq}[id=TOAS11-P2-PHY-S3-Q04, source={TOAS - Test 11 (27 Jun 2022)}, paper=2, qno=4, type=numerical, subtopic={Stefan-Boltzmann law}, answer={10.00}, answer_src=key, difficulty=5, diff_parts={K2 S2 I1 C0 F1}, diff_note={concepts: Stefan's law, Wien's law | steps: steady state with the source power; small temperature rise; new peak wavelength | I: linearise the temperature change}]
A spherical shell is kept in an atmosphere at temperature $T_0$. The wavelength corresponding to maximum intensity of radiation for the shell is $\lambda_0$. A point source of constant power is switched on inside the shell. The power radiated by the source is $P=0.4\,\sigma SeT_0^4$ where $S$, $e$ and $\sigma$ are outer surface area of the shell, emissivity of the outer surface of the shell and Stefan's constant respectively. Calculate the new wavelength ($\lambda$) corresponding to the maximum intensity of radiation from the shell. Assume that change in temperature ($\Delta T$) of the shell is small compared to the ambient temperature $T_0$. If value of $\lambda$ is given by $\left(\dfrac N{11}\right)\lambda_0$, then fill the value of $N$.
\end{bankq}
\begin{bankq}[id=TOAS19-P1-PHY-S3-Q05, source={TOAS - Test 19 (05 Feb 2023)}, paper=1, qno=5, type=numerical, subtopic={Radiation}, answer={3.00}, answer_src=key, difficulty=4, diff_parts={K1 S2 I1 C0 F1}, diff_note={concepts: radiation shields | steps: heat flux with each shield in series; steady state; ratio}]
A black surface at constant high temperature $T_H$ is parallel to another black plane surface at a constant lower temperature $T_L$ (enclosed volume in between). A heat shield consisting of two thin black plates is placed between the warm and the cold surfaces and parallel to these. After some time steady state is achieved. The ratio of initial heat current to final heat current is (Assume all surfaces to be large with respect to distance between surfaces)
\bfig{108.2}{TOAS19_p1_phy_s3_q05_1.png}
\end{bankq}

\section*{TAAS papers}
\begin{bankq}[id=TAAS02-P2-PHY-S4-Q07, source={TAAS - Test 02 (30 Jul 2023)}, paper=2, qno=7, type=numerical, subtopic={Stefan-Boltzmann law}, answer={9}, answer_src=key, difficulty=3, diff_parts={K1 S2 I0 C0 F1}, diff_note={concepts: Stefan-Boltzmann law | steps: temperature ratio in kelvin; power ratio; logarithm}]
A metal is heated in a furnace where a sensor is kept above the metal surface to read the power radiated ($P$) by the metal. The sensor has a scale that displays $\log_2(P/P_0)$, where $P_0$ is a constant. When the metal surface is at a temperature of $487^\circ$C, the sensor shows a value 1. Assume that the emissivity of the metallic surface remains constant. What is the value displayed by the sensor when the temperature of the metal surface is raised to $2767^\circ$C ?
\end{bankq}

\section*{TOPS papers}
\begin{bankq}[id=TOPS05-P1-PHY-S2-Q10, source={TOPS - Test 05 (14 Jan 2023)}, paper=1, qno=10, type=single, subtopic={Conduction}, answer={A}, answer_src=key, difficulty=5, diff_parts={K2 S2 I1 C1 F0}, diff_note={concepts: heat conduction, steady state | steps: initial conduction; steady rate constant; gradient and area; match}]
A copper rod (initially at room temperature of 20$^\circ$C) of non uniform cross section is placed between a steam chamber at 100$^\circ$C and Ice-water chamber at 0$^\circ$C.
\bfig{213.8}{TOPS05_p1_phy_s2_q10_1.png}
{\footnotesize\setlength{\tabcolsep}{3pt}
\begin{tabular}{|l|p{3.2cm}|l|p{2cm}|}
\hline
 & List--I & & List--II\\
\hline
(I) & Initially rate of heat flow $\left(\dfrac{dQ}{dt}\right)$ will be & (P) & Maximum at section $A$\\
(II) & At steady state rate of heat flow $\left(\dfrac{dQ}{dt}\right)$ will be & (Q) & Maximum at section $B$.\\
(III) & At steady state temperature gradient $\left|\dfrac{dT}{dx}\right|$ will be & (R) & Maximum at section $C$.\\
(IV) & At steady state rate of change of temperature $\left(\dfrac{dT}{dt}\right)$ at a certain point will be. & (S) & Minimum at section $B$.\\
 & & (T) & Same for all sections.\\
\hline
\end{tabular}}
\begin{qoptions}
\opt{A}{I$\to$P;II$\to$T;III$\to$Q;IV$\to$T}
\opt{B}{I$\to$P;II$\to$S;III$\to$Q;IV$\to$T}
\opt{C}{I$\to$P;II$\to$T;III$\to$T;IV$\to$S}
\opt{D}{I$\to$P;II$\to$T;III$\to$Q;IV$\to$S}
\end{qoptions}
\end{bankq}
\begin{bankq}[id=TOPS05-P2-PHY-S2-Q07, source={TOPS - Test 05 (14 Jan 2023)}, paper=2, qno=7, type=multi, subtopic={Stefan-Boltzmann law}, answer={A,C}, answer_src=key, difficulty=6, diff_parts={K2 S2 I1 C1 F1}, diff_note={concepts: Stefan law, Wien law | steps: area under the curve as T to the fourth power; temperature ratio; wavelength ratio}]
The spectral emission power $E_\lambda$ for a body at temperature $T_1$ and temperature $T_2$ is plotted against the wavelength and area under the curve is found $A$ for $T_1$ and $9A$ for $T_2$.
\bfig{116.3}{TOPS05_p2_phy_s2_q07_1.png}
\begin{qoptions}
\opt{A}{$\dfrac{T_1}{T_2}=\dfrac1{\sqrt3}$}
\opt{B}{$\dfrac{T_1}{T_2}=3$}
\opt{C}{$\dfrac{\lambda_1}{\lambda_2}=\sqrt3$}
\opt{D}{$\dfrac{\lambda_1}{\lambda_2}=\dfrac13$}
\end{qoptions}
\end{bankq}

\section*{TEAS papers}
\begin{bankq}[id=TEAS01-P1-PHY-S1-Q06, source={TEAS - Test 01 (16 Apr 2023)}, paper=1, qno=6, type=multi, subtopic={Stefan-Boltzmann law}, answer={B,C}, answer_src=key, difficulty=7, diff_parts={K2 S2 I2 C1 F1}, diff_note={concepts: Wien law, black body spectrum | steps: peak wavelength; shape of the spectrum; compare the intensities}]
The radiation coming from a black body at a temperature of 3000 K are analyzed. Assume wein's constant $b=3\times10^{-3}$ mK. We wish to compare intensity of radiation corresponding to wavelength range (999 nm to 1001 nm), say $E_1$; wavelength range (1499 nm to 1501 nm), say $E_2$; wavelength range (499 nm to 501 nm) say $E_3$.
\begin{qoptions}
\opt{A}{$E_2>E_1$}
\opt{B}{$E_2>E_3$}
\opt{C}{$E_1>E_3$}
\opt{D}{$E_3>E_1$}
\end{qoptions}
\end{bankq}

\section*{TMAS papers}
\begin{bankpassage}[id=TMAS03-P1-PHY-S2-G9-10, source={TMAS - Test 03 (10 Sep 2023)}, questions={9-10}]
A hollow cylinder has length $L$, inner radius $a$ and outer radius $b$. Thermal conductivity is $k$. Inner and outer surfaces of the cylinder are maintained at temperatures $100^\circ$C and $0^\circ$C respectively. If $kL=11$ J/s, $a=0.1$ m, $b=0.3$ m, and $\ell n3=1.1$.
\end{bankpassage}
\begin{bankq}[id=TMAS03-P1-PHY-S2-Q09, source={TMAS - Test 03 (10 Sep 2023)}, paper=1, qno=9, type=single, subtopic={Conduction}, answer={B}, answer_src=solved, working={$H=\dfrac{2\pi kL\,\Delta T}{\ln(b/a)}=\dfrac{2\pi\times11\times100}{1.1}=6.28$ kJ/s}, difficulty=3, diff_parts={K1 S1 I1 C1 F0}, diff_note={concepts: radial heat conduction | steps: formula for the cylinder; substitute}, group=TMAS03-P1-PHY-S2-G9-10]
The rate of radial heat (use $\pi=3.14$) flow (in kJ/s) is :-
\begin{qoptions}
\opt{A}{$8.26$}
\opt{B}{$6.28$}
\opt{C}{$5.75$}
\opt{D}{$7.65$}
\end{qoptions}
\end{bankq}
\begin{bankq}[id=TMAS03-P1-PHY-S2-Q10, source={TMAS - Test 03 (10 Sep 2023)}, paper=1, qno=10, type=single, subtopic={Conduction}, answer={A}, answer_src=solved, working={Gradient $=\dfrac{H}{2\pi kL\,a}=\dfrac{6280}{2\pi\times11\times0.1}=\dfrac{10000}{11}$, so $n=10$}, difficulty=3, diff_parts={K1 S1 I1 C1 F0}, diff_note={concepts: temperature gradient in a cylinder | steps: radial conduction law; gradient at the inner surface}, group=TMAS03-P1-PHY-S2-G9-10]
If the temperature gradient near the inner surface of cylinder is $\dfrac{1000n}{11}\ {}^\circ$C, then value of $n$ is :-
\begin{qoptions}
\opt{A}{$10$}
\opt{B}{$15$}
\opt{C}{$20$}
\opt{D}{$25$}
\end{qoptions}
\end{bankq}

\section*{TNAS papers}
\begin{bankq}[id=TNAS01-P1-PHY-S1-Q07, source={TNAS - Test 01 (05 Feb 2023)}, paper=1, qno=7, type=single, subtopic={Conduction}, answer={A}, answer_src=key, difficulty=5, diff_parts={K1 S2 I2 C1 F0}, diff_note={concepts: heat conduction in a varying cross section | steps: conduction law; integrate; temperature profile}]
A frustum shaped rod has cross sectional area that varies from its one end with distance $x$ as $A=A_0(1+\alpha x)$, where $A_0$ and $\alpha$ are constants. The rod has length $\ell$ and is maintained between the temperatures $T_1$ and $T_2$ ($T_1<T_2$) with $T_1$ being the temperature at cross section with smaller radius. Which of the following represents the temperature at a cross section at distance $x$ from the end at temperature $T_1$ ?
\begin{qoptions}
\opt{A}{$T_1+(T_2-T_1)\dfrac{\ell n\,(1+\alpha x)}{\ell n\,(1+\alpha\ell)}$}
\opt{B}{$T_1-T_2\dfrac{\ell n\,(1+\alpha x)}{\ell n\,(1+\alpha\ell)}$}
\opt{C}{$T_1-T_2\dfrac{\ell n\,(\alpha x)}{\ell n\,(\alpha\ell)}$}
\opt{D}{$T_2-T_1\dfrac{\ell n\,(1+\alpha x)}{\ell n\,(1+\alpha\ell)}$}
\end{qoptions}
\end{bankq}
\begin{bankq}[id=TNAS01-P2-PHY-S1-Q08, source={TNAS - Test 01 (05 Feb 2023)}, paper=2, qno=8, type=multi, subtopic={Radiation}, answer={A,B,C}, answer_src=key, difficulty=6, diff_parts={K1 S2 I2 C1 F1}, diff_note={concepts: radiation and conduction in a shell | steps: absorbed radiation; conduction formula; temperature relations}]
Imagine a point source $s$ which is a black body at a constant temperature $T$ at the centre of a thick spherical shell of inner radius $a$ and external radius $b$. The coefficient of thermal conductivity of the material of the shell is $K$. The inner surface of the shell is coated with lamp black. $A$ is the surface area of sources.
\bfig{82.0}{TNAS01_p2_phy_s1_q08_1.png}
\begin{qoptions}
\opt{A}{The amount of thermal radiation received by the inner surface from $s$ is ($\sigma\to$ Stefan's constant) $A\sigma T^4$.}
\opt{B}{The temperature difference across the inner and outer surface is $\dfrac{A\sigma T^4}{4\pi K}\left[\dfrac1a-\dfrac1b\right]$.}
\opt{C}{If the external surface of the shell is in vacuum at $0$ K (say) the temperature of inner surface, is $\dfrac{A\sigma T^4}{4\pi K}\left[\dfrac1a-\dfrac1b\right]$.}
\opt{D}{None of these}
\end{qoptions}
\end{bankq}
\begin{bankq}[id=TNAS04-P1-PHY-S1-Q02, source={TNAS - Test 04 (18 Jun 2023)}, paper=1, qno=2, type=single, subtopic={Conduction}, answer={D}, answer_src=key, difficulty=4, diff_parts={K1 S2 I1 C1 F0}, diff_note={concepts: heat conduction with variable conductivity | steps: constant heat current; separate variables; integrate to get temperature profile}]
The ends of a rod of length $\ell$ are maintained at $\theta_1$ and $\theta_2$ ($<\theta_1$). The thermal conductivity of rod at a point is given as $k=\dfrac{\alpha}{\theta}$ where $\theta$ = temperature at that point of the rod $\theta$ as a function of $x$ is.
\bfig{147.5}{TNAS04_p1_phy_s1_q02_1.png}
\begin{qoptions}
\opt{A}{$\dfrac{x}{\ell}\dfrac{\theta_1^2}{\theta_2}$}
\opt{B}{$\dfrac{\theta_1 x}{\ell}e^{\theta_1/\theta_2}$}
\opt{C}{$\left(\dfrac{x}{\ell}\right)^2\dfrac{\theta_1^2}{\theta_2^2}$}
\opt{D}{$\theta_1\left(\dfrac{\theta_2}{\theta_1}\right)^{x/\ell}$}
\end{qoptions}
\end{bankq}
\begin{bankq}[id=TNAS04-P1-PHY-S2-Q11, source={TNAS - Test 04 (18 Jun 2023)}, paper=1, qno=11, type=multi, subtopic={Conduction}, answer={A,B,C,D}, answer_src=key, difficulty=9, diff_parts={K2 S3 I2 C2 F1}, diff_note={concepts: heat exchanger with counterflow, heat conduction through a wall | steps: heat flow across wall; fluid temperature change; same difference everywhere; energy balance; final difference}]
A heat exchanger is a device which transfers heat between a hot fluid and a cold fluid. The heat exchanger shown below consists of two rectangular tubes of length $\ell$, width $w$ and height $h$. The tubes are separated by a metal wall of thickness $d$ and thermal conductivity $k$. Originally hot fluid flows through the lower tube at a speed $v$ from right to left, and originally cold fluid flows through the upper tube in the opposite direction (left to right) at the same speed. The heat capacity per unit volume of both fluids is $c$.\par
The hot fluid enters the heat exchanger at a higher temperature than the cold fluid; the difference between the temperatures of the entering fluids is $\Delta T_i$. When the fluids exit the heat exchanger the difference has been reduced to $\Delta T_f$.
\bfig{209.7}{TNAS04_p1_phy_s2_q11_1.png}
Assume that the temperature in each pipe changes linearly with the length of the tube, and consider transfer of heat only due to conduction in the metal and due to the bulk movement of fluid. Under the assumptions in this problem, while the temperature of each fluid varies along the length of the exchanger, the temperature difference across the wall is the same everywhere.
\begin{qoptions}
\opt{A}{The amount of energy transferred from one liquid to another per unit time is $\dfrac{chwv\Delta T_i}{\left(1+\frac{chvd}{\ell k}\right)}$.}
\opt{B}{The change in temperature of a liquid as it crosses the length $\ell$ is $\dfrac{\Delta T_i}{\left(1+\frac{chvd}{\ell k}\right)}$.}
\opt{C}{The increase in temperature of one liquid is equal to decrease in temperature of other liquid.}
\opt{D}{The value of $\Delta T_f$ is $\left[1-\dfrac{2}{1+\frac{chvd}{\ell k}}\right]\Delta T_i$}
\end{qoptions}
\end{bankq}
\begin{bankq}[id=TNAS04-P1-PHY-S3-Q01, source={TNAS - Test 04 (18 Jun 2023)}, paper=1, qno=1, type=single, subtopic={Conduction}, answer={3}, answer_src=solved, working={Resistance of the shell $R=\int\dfrac{r\,dr}{4\pi C}=\dfrac{3a^2}{8\pi C}$; $mS\,\dfrac{d\theta}{dt}=-\dfrac{\theta-30}{R}$ gives $t=RmS\ln\dfrac{30}{10}=3\,\dfrac{mSa^2\ln3}{8\pi C}$, so $n=3$ (the printed $a^3$ is a dimensional misprint for $a^2$)}, difficulty=5, diff_parts={K2 S2 I1 C1 F0}, diff_note={concepts: conduction through a spherical shell, cooling of water | steps: thermal resistance by integration; rate equation for water; integrate from sixty to forty; read off n}]
Water of mass $m$ and specific heat capacity $S$ is kept at $60^\circ$ C inside a thick spherical shell of inner and outer radii $a$ and $2a$ as shown. Thermal conductivity of the shell varies as $k=\dfrac{C}{r^3}$, where $C$ is a constant and $r$ is the distance from the centre. The surrounding temperature is $\theta_0=30^\circ$C. If the time required to decrease the temperature of water from $60^\circ$C to $40^\circ$C is $t=n\left(\dfrac{mSa^3\ell n3}{8\pi C}\right)$. Find the value of $n$.
\bfig{71.5}{TNAS04_p1_phy_s3_q01_1.png}
\end{bankq}
\begin{bankq}[id=TNAS04-P2-PHY-S1-Q05, source={TNAS - Test 04 (18 Jun 2023)}, paper=2, qno=5, type=multi, subtopic={Conduction}, answer={A,D}, answer_src=key, difficulty=9, diff_parts={K2 S3 I2 C2 F1}, diff_note={concepts: conduction through a freezing layer, spherical geometry | steps: thermal resistance of the ice shell; rate of freezing; total time by integration; time to freeze seven eighths; compare with the last eighth}]
A thin hollow spherical shell of radius $R$ is filled with water. The temperature of atmosphere is $T_s$ which is below $0^\circ$C. Firstly all water comes down to $0^\circ$C and then starts freezing from the outer surface towards the centre. Assume that the thermal conductivity of shell is very high. Neglect change in density due to freezing. $\rho$ = density of ice, $k$ = thermal conductivity of ice, $L$ = latent heat of fusion.
\begin{qoptions}
\opt{A}{Time taken to whole water to freeze is $\dfrac{\rho LR^2}{6kT_S}$}
\opt{B}{Time taken to whole water to freeze is $\dfrac{2\rho LR^2}{kT_S}$.}
\opt{C}{Time taken to freeze first $\dfrac78$th of total mass is more than time taken for last $\dfrac M8$ mass to freeze.}
\opt{D}{Time taken to freeze first $\dfrac78$th of total mass is same as time taken for last $\dfrac M8$ mass to freeze.}
\end{qoptions}
\end{bankq}
\begin{bankq}[id=TNAS04-P2-PHY-S3-Q04, source={TNAS - Test 04 (18 Jun 2023)}, paper=2, qno=4, type=numerical, subtopic={Newton's law of cooling}, answer={910.00}, answer_src=key, difficulty=6, diff_parts={K2 S2 I1 C1 F1}, diff_note={concepts: Newton's law of cooling, water equivalent | steps: cooling rate gives the loss; equal to the power input; evaluate}]
A heating element is used to raise the temperature of water of mass 10 kg inside the boiler. The heater is kept on for a very long time and it is observed that no steam is formed. Now the heating element is switched off and temperature of water is observed to decrease at a rate of $1^\circ$C/min. The water equivalent of boiler is 3 kg. The power (in W) rating of heating element is (Take specific heat capacity of water 4200 J/kg-K)
\end{bankq}

\section*{Internal Test papers}
\begin{bankq}[id=INT03-P1-PHY-S1-Q04, source={Internal Test - Test 03 (08 Sep 2021)}, paper=1, qno=4, type=single, subtopic={Conduction}, answer={B}, answer_src=key, difficulty=5, diff_parts={K1 S2 I2 C1 F0}, diff_note={concepts: thickness of ice formed on a lake | steps: steady conduction through ice and water; equal heat currents; solve for the thickness}]
In a 20 m deep lake, the bottom is at a constant temperature of $4^\circ$C. The air temperature is constant at $-10^\circ$C. The thermal conductivity of ice is 4 times that water. Neglecting the expansion of water on freezing, the maximum thickness of ice will be
\begin{qoptions}
\opt{A}{$\dfrac{20}{11}$ m}
\opt{B}{$\dfrac{200}{11}$ m}
\opt{C}{$20$ m}
\opt{D}{$10$ m}
\end{qoptions}
\end{bankq}
\begin{bankq}[id=INT03-P1-PHY-S2-Q12, source={Internal Test - Test 03 (08 Sep 2021)}, paper=1, qno=12, type=multi, subtopic={Radiation}, answer={B,D}, answer_src=key, difficulty=5, diff_parts={K1 S2 I2 C0 F1}, diff_note={concepts: cooling of a sphere and a cube of equal surface area | steps: rate of heat loss from the surface; heat content from volume; compare rate of fall}]
A hollow copper sphere \& a hollow copper cube of same surface area \& negligible thickness, are filled with warm water of same temperature and placed in an enclosure of constant temperature, a few degrees below that of the bodies. Then in the beginning :
\begin{qoptions}
\opt{A}{The rate of energy lost by the sphere is greater than that by the cube}
\opt{B}{The rate of energy lost by the two are equal}
\opt{C}{The rate of energy lost by the sphere is less than that by the cube}
\opt{D}{The rate of fall of temperature for sphere is less than that for the cube.}
\end{qoptions}
\end{bankq}
\begin{bankq}[id=INT03-P1-PHY-S3-Q16, source={Internal Test - Test 03 (08 Sep 2021)}, paper=1, qno=16, type=numerical, subtopic={Conduction}, answer={2}, answer_src=key, difficulty=6, diff_parts={K1 S2 I2 C1 F1}, diff_note={concepts: ratio of thermal conductivities from temperature profiles | steps: slopes of the graphs; equal heat flow; ratio}]
Two different rods A and B are kept as shown in figure.\par
\bfig{207.3}{INT03_p1_phy_s2_q16_1.png}
\bfig{163.2}{INT03_p1_phy_s2_q16_2.png}
The variation of temperature of different cross sections is plotted in a graph shown in figure. Find the ratio of thermal conductivities of B to A.
\end{bankq}
\begin{bankq}[id=INT03-P1-PHY-S3-Q17, source={Internal Test - Test 03 (08 Sep 2021)}, paper=1, qno=17, type=numerical, subtopic={Newton's law of cooling}, answer={2}, answer_src=key, difficulty=6, diff_parts={K1 S2 I2 C1 F1}, diff_note={concepts: Newton's law of cooling for different temperatures | steps: rate of cooling is proportional to the excess; ratio of the excess temperatures; time}]
A body cools from $50^\circ$C to $49.9^\circ$C in 5 sec. It cools from $40^\circ$C to $39.9^\circ$C in $t$ sec. Assuming Newtons law of cooling to be valid and temperature of surrounds at $30^\circ$C, value of $t/5$ will be? (in sec) (to the nearest integer)
\end{bankq}
\begin{bankq}[id=INT03-P1-PHY-S3-Q18, source={Internal Test - Test 03 (08 Sep 2021)}, paper=1, qno=18, type=numerical, subtopic={Conduction}, answer={9}, answer_src=key, difficulty=8, diff_parts={K1 S3 I2 C2 F1}, diff_note={concepts: melting of ice and evaporation of water through a conducting rod | steps: heat flow into each end; latent heat relation; temperature gradients; geometry; solve for the distance}]
A metal rod AB of length $10x$ has its one end A in ice at $0^\circ$C and the other end B in water at $100^\circ$C. It a point P on the rod is maintained at $400^\circ$C, then it is found that equal amounts of water and ice evaporate and melt per unit time respectively. The latent heat of evaporation of water 540 cal/g and latent heat of melting of ice is 80 cal/g. If the point P is at $lx$ distance from the ice end then find out the value of $l$. [Neglect any heat loss to the surrounding.]
\end{bankq}
\end{document}
